% ============================================================
\begin{abstract}
Stratospheric sudden warming (SSW) events produce persistent surface weather anomalies, yet their consequences for geophysical hazards remain unexplored.
Here we identify a previously unknown connection: SSW episodes drive a large, reproducible reduction in natural dry slab avalanche activity across four countries and two continents, mediated by a ``threshold amplification'' mechanism in which modest stratospheric-origin weather shifts produce disproportionately large hazard responses by switching binary trigger conditions across mountain networks.
In Switzerland (21~winters, 16~SSW events), 14 of 16 events show reduced counts (geometric mean rate ratio $\mathrm{RR} = 0.32$; 95\% CI $[0.20, 0.54]$; $d = -1.06$).
Norwegian avalanche danger levels decrease concordantly ($d = -0.67$; $P < 10^{-6}$), and all four Utah SSW events show decreased counts ($\mathrm{RR} = 0.34$).
A specification curve of 180 analytical variants confirms robustness: 100\% show $\mathrm{RR} < 1$ (permutation $P < 0.001$).
Five-country EAWS data reveal a predictable geographic gradient in hazard response ($r = 0.69$, $P = 0.0004$), confirmed out-of-sample by independent ALBINA bulletin data---demonstrating that the spatial structure of SSW--surface coupling is systematic, not stochastic.
A count--rating dissociation---danger \emph{ratings} increase while occurrence \emph{counts} decrease---and a differential trigger response---human-triggered avalanches increase ($\mathrm{RR} = 1.81\times$ relative to natural) while natural triggers are suppressed---reveal a ``loaded gun'' mechanism: SSW-associated cold regimes suppress surface-energy triggers ($-57\%$ warming, $-29\%$ rain) while preserving snowpack instability.
Hemispheric verification via Aura/MLS stratospheric chemistry, SNOTEL surface observations, and Northern Hemisphere snow cover ($d = +1.13$) confirms that every link in the mechanism chain is independently significant.
The temporal profile---suppression emerging before SSW onset---identifies planetary wave forcing as the common cause rather than top-down stratospheric forcing.
\end{abstract}

\noindent\textbf{Plain Language Summary.}
Snow avalanches in the Alps kill hundreds of people each year. We discovered that a dramatic atmospheric event called a ``stratospheric sudden warming''---where temperatures 30--50~km above Earth surge by 30--40~°C in just a few days---is consistently followed by a sharp drop in natural avalanche activity at the surface. Across 16 such events in Switzerland over 21~winters, natural dry slab avalanches decreased by 68\% on average. The same pattern appears in Norway and Utah using entirely different measurement systems. Using output from SNOWPACK, a physics-based snowpack model run at 130~Swiss stations, we find that the snowpack actually becomes \emph{less} stable during these events---but the natural triggers that would normally release avalanches (surface warming, rain-on-snow, solar radiation) are simultaneously suppressed by the cold regime. This creates a ``loaded gun'' situation: a primed snowpack without natural release, which may explain why human-triggered avalanche rates increase during these windows. Because stratospheric sudden warmings can be predicted 1--2~weeks in advance, our findings could extend avalanche hazard forecasting beyond the current 5--7~day weather-based window.

